\section{Local options of the macro  \tkzcname{AQquestion}}

\subsection{Local use of \tkzname{pq}}
\Iopt{AQquestion}{pq} 
 The following table is obtained with the options |lq=85mm| and |size=\wide|. The questions are misplaced. The local option \tkzname{pq} solves this problem, the text can be moved 1mm upwards with \tkzcname{AQquestion[pq=1mm]}. 
  and by |6mm| for the second.

\medskip 


 \begin{alterqcm}[lq=55mm,size=\large]

\AQquestion{If the function $f$ is strictly increasing on $\mathbf{R}$ then the equation $f(x) = 0$ admits :}
{{At least one solution},
[At most one solution],
{Exactly one solution}
}
\AQquestion{If the $f$ function is continuous and positive on $[a~ ;~ b]$ and $\mathcal{C}_{f}$ its representative curve in an orthogonal system. In units of area, the area $\mathcal{A}$ of the domain delimited by $\mathcal{C}_{f}$, the abscissa axis and the lines of equations $x = a$ 5 and $x = b$ is given by the formula : }
{%
{$\mathcal{A}= \displaystyle \int_{b}^a f(x)\ \text{d}x$},
{$\mathcal{A}= \displaystyle \int_{a}^b f(x)\ \text{d}x$},
{$\mathcal{A} = f(b) - f(a)$}}
\end{alterqcm}

\medskip 
\tkzname{Here is the corrected version}

\begin{alterqcm}[lq=55mm,size=\large]
\AQquestion[pq=1mm]{If the $f$ function is strictly increasing on
$\mathbf{R}$ then the equation $f(x) = 0$ admits...}
{{At least one solution},
{At most one solution},
{Exactly one solution}
}
\AQquestion[pq=6mm]{If the $f$ function is continuous and positive on $[a~ ;~ b]$
 and $\mathcal{C}_{f}$ its representative curve in an orthogonal system.
  In area units, the $\mathcal{A}$ area of the domain delimited by $\mathcal{C}_{f}$, the abscissa axis and the lines of equations $x = a$ and $x = b$ is given  by the formula: }
{{$\mathcal{A}= \displaystyle \int_{b}^a f(x)\ \text{d}x$},
{$\mathcal{A}= \displaystyle \int_{a}^b f(x)\ \text{d}x$},
{$\mathcal{A} = f(b) - f(a)$}
}
\end{alterqcm}

\medskip
\begin{tkzexample}[code only, small]
 \begin{alterqcm}[lq=55mm,size=\large]
 \AQquestion[pq=1mm]{If the $f$ function is strictly increasing on  $\mathbf{R}$
 then the equation $f(x) =0 $ admits...
 {{{At least one solution},
 [At most one solution],
 {Exactly one solution}}
\end{tkzexample}

\medskip
\begin{tkzexample}[code only, small]
 \AQquestion[pq=6mm]{If the $f$ function is continuous and positive on $[a~ ;~ b]$ and $\mathcal{C}_{f}$ its representative curve in an orthogonal system.
   In units of area, the area $\mathcal{A}$ of the domain delimited by $\mathcal{C}_{f}$, the abscissa axis and the lines of equations $x = a$ and $x = b$ is given by the formula: }
 {{$\mathcal{A}= \displaystyle \int_{b}^a f(x)\ \text{d}x$},
 {$\mathcal{A}= \displaystyle \int_{a}^b f(x)\ \text{d}x$},
 {$\mathcal{A} = f(b) - f(a)$}}
 \end{alterqcm}
\end{tkzexample}

\subsection{Global and local use of \tkzname{pq}}\
 \Iopt{AQquestion}{pq} \IoptEnv{alterqcm}{pq}
This time, it is necessary to move several questions, I placed a |pq=2mm| globally, that is to say like this :\tkzcname{begin\{alterqcm\}[lq=85mm,pq=2mm]}. \textbf{All} questions are affected by this option but some questions were well placed and should remain so, so locally I give them back a |pq=0mm|.

\medskip
\begin{alterqcm}[lq=85mm,pq=2mm]
\AQquestion{A bivariate statistical series. The values of $x$ are 1, 2, 5, 7, 11, 13 and a least squares regression line equation of $y$ to $x$ is $y = 1.35x +22.8$. The coordinates of the mean point are :}
{{$(6,5;30,575)$},
{$(32,575 ; 6,5)$},
{$(6,5 ; 31,575)$}}

\AQquestion[pq=0mm]{$(u_{n})$ is an arithmetic sequence of reason $-5$.\\ Which of these statements is true? }
{{For all $n,~  u_{n+1} - u_{n} = 5$},
{$u_{10}= u_{2}+ 40$},
{$u_{3} = u_{7} + 20$}
}
\AQquestion[pq=0mm]{Equality $\ln (x^2 - 1) = \ln (x - 1) + \ln (x+1)$ is true}
{{For all $x$ in  $]- \infty~;~-1[ \cup]1~;~+ \infty[$},
{For all $x$ in $\mathbf{R} - \{-1~ ;~ 1\}$.},
{For all $x$ in $]1~ ;~+\infty[$}
}
\AQquestion{For all $x$, the number \[\dfrac{\text{e}^x - 1}{\text{e}^x + 2}\hskip12pt \text{equal to :} \] }
{{$-\dfrac{1}{2}$},
{$\dfrac{\text{e}^{-x} - 1}{\text{e}^{-x} + 2}$},
{$\dfrac{1 - \text{e}^{-x}}{1 + 2\text{e}^{-x}}$}
}
\AQquestion{Let I $= \displaystyle\int_{\ln 2}^{\ln 3} \dfrac{1}{\text{e}^x - 1}\,\text{d}x$ and J $ = \displaystyle\int_{\ln 2}^{\ln 3} \dfrac{\text{e}^x}{\text{e}^x - 1}\,\text{d}x$ \\ then the number  I $-$ J is equal to}
{{$\ln \dfrac{2}{3}$},
{$\ln \dfrac{3}{2}$},
{$\dfrac{3}{2}$}
}
\end{alterqcm}

\medskip
\begin{tkzexample}[code only,vbox,small]
 \begin{alterqcm}[lq=85mm,pq=2mm]
  \AQquestion[pq=0mm]{Equality $\ln (x^2 - 1) = \ln (x - 1) + \ln (x+1)$
   is true}
  {{For all $x$ in  $]- \infty~;~-1[ \cup]1~;~+ \infty[$},
  {For all $x$ in $\mathbf{R} - \{-1~ ;~ 1\}$.},
  {For all $x$ in $]1~ ;~+\infty[$}}
  \AQquestion{For any real $x$, the number \[\dfrac{\text{e}^x - 1}
  {\text{e}^x + 2}\hskip12pt \text{equal to :} \] }
  {{$-\dfrac{1}{2}$},
  {$\dfrac{\text{e}^{-x} - 1}{\text{e}^{-x} + 2}$},
  {$\dfrac{1 - \text{e}^{-x}}{1 + 2\text{e}^{-x}}$}}
  \end{alterqcm}
 \end{tkzexample}


\subsection{\tkzname{correction} and \tkzname{br} : rank of good answer}
\Iopt{AQquestion}{br}  \Iopt{AQquestion}{correction}
First of all, it is necessary to ask for an answer key. To do this, just include the option \tkzname{correction} which is a boolean, thus set to \tkzname{true}. Then in each question, it is necessary to give the list of correct answers. For example, with \tkzname{br=1} or \tkzname{br=\{1,3\}}.

Here is the previous year's correction:

\medskip
\begin{tkzexample}[vbox,small]
\begin{alterqcm}[VF,correction,lq=125mm]
 \AQquestion[br=1]{For all $x \in ]-3~;~2],~f'(x) \geqslant 0$.}
 \AQquestion[br=2]{The $F$ function has a maximum in $2$}
 \AQquestion[br=2]{$\displaystyle\int_{0}^2 f'(x)\:\text{d}x = - 2$}
\end{alterqcm}
\end{tkzexample}


\endinput